% ch03_u03 —— 原书 p.81–p.90（3.5 中段至 3.6 开头）
\SourcePage{81}
3.5 与 3.6 两节的思路\footnote{受 [Phys] 启发。}是：给 $A_{\mathfrak{g}}(O)$ 及其对应物 $\mathfrak{z}_{{}^{L}\mathfrak{g}}(O)$ 附加某些结构，以「消去」它们的自同构。对 $A_{\mathfrak{g}}(O)$ 而言，所附加的是 3.5.11 中的李代数胚 $\mathfrak{a}_{\mathfrak{g}}$。它在 $\mathfrak{z}_{{}^{L}\mathfrak{g}}(O)$ 中的对应物将在 3.6.5 中引入。$\mathfrak{a}_{\mathfrak{g}}$ 的定义很简单：它是 $\operatorname{Op}_{\mathfrak{g}}(O)$ 上万有 $G$-丛 $\mathfrak{F}^{0}_{G}$ 的无穷小对称代数胚。$\mathfrak{F}^{0}_{G}$ 从而 $\mathfrak{a}_{\mathfrak{g}}$ 都带有 $\operatorname{Der}O$ 的作用。结果是，偶 $(A_{\mathfrak{g}}(O),\mathfrak{F}^{0}_{G})$ 没有非平凡的 $\operatorname{Der}O$-等变自同构（见 3.5.9）；对 $(A_{\mathfrak{g}}(O),\mathfrak{a}_{\mathfrak{g}})$ 而言，这一点「几乎」成立（见 3.5.13）。

\bdp{3.5.1} 我们有 $\operatorname{Spec}O$ 上由 3.2.1 中的概形 $\operatorname{Op}_{\mathfrak{g}}(O)=\operatorname{Spec}A_{\mathfrak{g}}(O)$ 参数化的通用 $\mathfrak{g}$-oper 族。固定一个一维自由 $O$-模 $\omega_{O}^{1/2}$，并配上同构 $\omega_{O}^{1/2}\otimes\omega_{O}^{1/2}\xrightarrow{\sim}\omega_{O}$（当然，$\omega_{O}^{1/2}$ 在同构意义下唯一）。于是上述通用族提升为 $G$-oper 族；见 3.4.1\footnote{说实话，我们还须对每个负单根 $\alpha$ 选取非零的 $y_{\alpha}\in\mathfrak{g}^{\alpha}$（见 3.4.1）。}。这样，我们得到 $\operatorname{Spec}(A_{\mathfrak{g}}(O)\widehat\otimes O)=\operatorname{Spec}A_{\mathfrak{g}}(O)[[t]]$ 上的 $B$-丛 $\mathfrak{F}_{B}$，以及相伴 $G$-丛 $\mathfrak{F}_{G}$ 上沿 $\operatorname{Spec}O$ 的联络 $\nabla$。

\bdp{3.5.2} 考虑群归纳概形 $\operatorname{Aut}_{2}O:=\operatorname{Aut}(O,\omega_{O}^{1/2})$。我们有典范正合列
\begin{equation}\tag{63}
0\to\mu_{2}\to\operatorname{Aut}_{2}O\to\operatorname{Aut}O\to0,
\end{equation}
且 $\operatorname{Aut}_{2}O$ 连通。正合列 (63) 与连通性可以作为 $\operatorname{Aut}_{2}O$ 的另一定义。以 $\operatorname{Aut}_{2}^{0}O$ 记 $\operatorname{Aut}^{0}O$ 在 $\operatorname{Aut}_{2}O$ 中的原像。

$\operatorname{Aut}O$ 作用于 $A_{\mathfrak{g}}(O)$ 与 $O$，从而作用于 $\operatorname{Spec}(A_{\mathfrak{g}}(O)\widehat\otimes O)$。该作用典范地提升为 $\operatorname{Aut}_{2}O$ 在 $(\mathfrak{F}_{B},\nabla)$ 上的作用。$\mu_{2}\subset\operatorname{Aut}_{2}O$ 经态射 (56) 作用于 $\mathfrak{F}_{B}$。

\SourcePage{82}
\bdp{3.5.3} \emph{引理.} 设 $L$ 是代数群，$A$ 是带 $\operatorname{Aut}O$ 作用的代数。考虑 $\operatorname{Aut}O$ 经其在 $A$ 与 $O$ 上的作用诱导的、在 $A\widehat\otimes O$ 上的作用。令 $i:\operatorname{Spec}A\hookrightarrow\operatorname{Spec}(A\widehat\otimes O)$ 为自然嵌入，$\pi:\operatorname{Spec}(A\widehat\otimes O)\to\operatorname{Spec}A$ 为投影。

\begin{enumerate}[label=\arabic*),leftmargin=2.2em]
\item $i^{*}$ 是 $\operatorname{Spec}(A\widehat\otimes O)$ 上 $\operatorname{Aut}_{2}O$-等变 $L$-丛的范畴与 $\operatorname{Spec}A$ 上 $\operatorname{Aut}_{2}^{0}O$-等变 $L$-丛的范畴之间的等价。
\item $\pi^{*}$ 是 $\operatorname{Spec}A$ 上 $\operatorname{Aut}_{2}O$-等变 $L$-丛的范畴与 $\operatorname{Spec}(A\widehat\otimes O)$ 上带沿 $\operatorname{Spec}O$ 的 $\operatorname{Aut}_{2}O$-不变联络的 $\operatorname{Aut}_{2}O$-等变 $L$-丛的范畴之间的等价。
\item 这些等价与遗忘函子（从 $\operatorname{Spec}A$ 上 $\operatorname{Aut}_{2}O$-等变丛到 $\operatorname{Spec}A$ 上 $\operatorname{Aut}_{2}^{0}O$-等变丛）及遗忘函子（从带联络的丛到丛）相容。\hfill$\Box$
\end{enumerate}

\bdp{3.5.4} 以 $\mathfrak{F}^{0}_{B}$ 与 $\mathfrak{F}^{0}_{G}$ 记 $\mathfrak{F}_{B}$ 与 $\mathfrak{F}_{G}$ 到 $\operatorname{Op}_{\mathfrak{g}}(O)=\operatorname{Spec}A_{\mathfrak{g}}(O)\subset\operatorname{Spec}A_{\mathfrak{g}}(O)\widehat\otimes O$ 的限制。$\mathfrak{F}^{0}_{B}$ 是 $\operatorname{Op}_{\mathfrak{g}}(O)$ 上的 $B$-丛，$\mathfrak{F}^{0}_{G}$ 是相应的 $G$-丛。$\mathfrak{F}^{0}_{B}$ 是 $\operatorname{Aut}_{2}^{0}O$-等变的；由 3.5.3，$\mathfrak{F}^{0}_{G}$ 是 $\operatorname{Aut}_{2}O$-等变的。由于 $\mathfrak{F}_{G}$ 上的联络 $\nabla$ 不保持 $\mathfrak{F}_{B}$，$\operatorname{Aut}_{2}O$ 在 $\mathfrak{F}^{0}_{G}$ 上的作用不保持 $\mathfrak{F}^{0}_{B}$。由 3.5.3，带 $\operatorname{Aut}_{2}O$ 作用的 $\mathfrak{F}^{0}_{G}$ 与 $B$-结构 $\mathfrak{F}^{0}_{B}\subset\mathfrak{F}^{0}_{G}$ 一起「记得」通用 oper $(\mathfrak{F}_{B},\nabla)$。

\bdp{3.5.5} 以 $F^{0}_{H}$ 记 $\operatorname{Op}_{\mathfrak{g}}(O)$ 上对应于 $\mathfrak{F}^{0}_{B}$ 的 $H$-丛。由于 $\operatorname{Op}_{\mathfrak{g}}(O)$ 是（无穷维）仿射空间，$\operatorname{Op}_{\mathfrak{g}}(O)$ 上任何 $H$-丛都平凡，且 $H$ 在其平凡化集合上的作用传递。特别地，这适用于 $\mathfrak{F}^{0}_{H}$；于是 $\mathfrak{F}_{H}$ 是 $\operatorname{Spec}\mathbb{C}$ 上某个 $H$-丛 $F_{H}$ 的拉回。由 3.4.1，$F_{H}$ 是 $\operatorname{Spec}\mathbb{C}$ 上的 $\mathbb{G}_{m}$-丛 $\omega_{O}^{1/2}/t\omega_{O}^{1/2}$ 沿 (54) 所定义态射 $\lambda^{\#}:\mathbb{G}_{m}\to H$ 的前推。特别地，$\operatorname{Aut}_{2}^{0}O$ 在 $F_{H}$ 上的作用来自复合
\[
\operatorname{Aut}_{2}^{0}O\to\operatorname{Aut}(\omega_{O}^{1/2}/t\omega_{O}^{1/2})=\mathbb{G}_{m}\xrightarrow{\lambda^{\#}}H.
\]

\SourcePage{83}
于是 $\operatorname{Der}^{0}O$ 在 $F_{H}$ 上的作用是「显然」作用（即保持 $F_{H}$ 的任一平凡化的作用）与态射 $\operatorname{Der}^{0}O\to\mathfrak{h}$ 之和，后者由 $f(t)\cdot t\frac{d}{dt}\mapsto f(0)\check\rho$ 定义。这里 $\check\rho$ 是基本余权之和。

\bdp{3.5.6} 这里给出 $A_{\mathfrak{g}}(O)$ 与 $\mathfrak{F}^{0}_{G}$ 的一个显式描述，其精神与 3.1.9–3.1.10 相同。设 $e,f\in\mathfrak{sl}_{2}$ 为 3.1.8 中的矩阵。固定主嵌入 $i:\mathfrak{sl}_{2}\hookrightarrow\mathfrak{g}$，使得 $i(e)\in\mathfrak{b}$。若选取嘉当子代数 $\mathfrak{h}\subset\mathfrak{b}$ 使 $i([e,f])\in\mathfrak{h}$，则 $i([e,f])$ 可与 $2\check\rho$ 等同。正如 3.1.9，置 $V:=\operatorname{Ker}\operatorname{ad}i(e)$。选取基 $e_{1},\dots,e_{r}\in V$，使得 $e_{1}=i(e)$，且所有 $e_{j}$ 都是 $\operatorname{ad}\check\rho$ 的特征向量。事实上，$[\check\rho,e_{j}]=(d_{j}-1)e_{j}$，其中 $d_{j}$ 是 $\mathfrak{g}$ 上「基本」不变多项式的次数（特别地 $d_{1}=2$）。平凡 $G$-丛上的联络
\begin{equation}\tag{64}
\nabla_{\frac{d}{dt}}=\frac{d}{dt}+i(f)+u_{1}(t)e_{1}+\cdots+u_{r}(t)e_{r}
\end{equation}
定义一个 $\mathfrak{g}$-oper；由 3.1.10，这在 $\operatorname{Spec}O$ 上的 $\mathfrak{g}$-oper（$O:=\mathbb{C}[[t]]$）与 $r$-元组 $(u_{1}(t),\dots,u_{r}(t))$、$u_{j}(t)\in\mathbb{C}[[t]]$ 之间给出双射。把 $u_{j}(t)$ 写成 $u_{j0}+u_{j1}t+\cdots$。于是 $A_{\mathfrak{g}}(O)$ 是 $u_{jk}$（$1\le j\le r$，$0\le k<\infty$）的多项式环。3.5.1 与 3.5.4 中的丛 $\mathfrak{F}_{B},\mathfrak{F}_{G},\mathfrak{F}^{0}_{B},\mathfrak{F}^{0}_{G}$ 都平凡；我们通过选取 opers 的典范形式 (64) 将它们平凡化。

为描述 $\operatorname{Der}O$ 在 $A_{\mathfrak{g}}(O)$ 与 $\mathfrak{F}^{0}_{G}$ 上的作用，引入标准记号 $L_{n}:=-t^{n+1}\frac{d}{dt}\in\operatorname{Der}\mathbb{C}((t))$（于是 $n\ge-1$ 时 $L_{n}\in\operatorname{Der}O$）。置 $u_{j}:=u_{j0}$。则
\begin{equation}\tag{65}
u_{jk}=(L_{-1})^{k}u_{j}/k!
\end{equation}
\begin{equation}\tag{66}
L_{0}u_{j}=d_{j}u_{j}
\end{equation}
\begin{equation}\tag{67}
L_{n}u_{j}=0\qquad\text{若 }n>0,\ j\ne1
\end{equation}
\begin{equation}\tag{68}
L_{n}u_{1}=0\qquad\text{若 }n>0,\ n\ne2;\qquad L_{2}u_{1}=-3.
\end{equation}
于是 $A_{\mathfrak{g}}(O)$ 是以 $u_{1},\dots,u_{r}$ 为生成元、以关系 (66)–(68) 为定义关系的交换 $(\operatorname{Der}O)$-代数。以 $L_{n}^{hor}$ 记 $\mathfrak{F}^{0}_{G}$ 上的向量场，它\SourcePage{84}%
来自 $\mathfrak{F}^{0}_{G}$ 的平凡化与 $\operatorname{Der}O$ 在 $A_{\mathfrak{g}}(O)$ 上的作用。

$L_{n}$ 在 $\mathfrak{F}^{0}_{G}$ 上的作用为 $L_{n}^{hor}+M_{n}$，其中 $M_{n}\in\mathfrak{g}\otimes A_{\mathfrak{g}}(O)$。可以证明
\begin{equation}\tag{69}
M_{0}=-\check\rho
\end{equation}
\begin{equation}\tag{70}
M_{1}=-i(e),\qquad M_{n}=0\quad(n>1)
\end{equation}
\[
M_{-1}=i(f)+u_{1}e_{1}+\cdots+u_{r}e_{r}.
\]
后文只会用到 (69)（恐怕至少还会用到 (70)）!???!）。
% [Tutor Note: 原书 p.84 此处作“... (I’m afraid we’ll use at least (70)) !???!).”，括号不配对，底本如此，保留原样。]

\emph{注.} 若 $n\ge0$，则 $M_{n}\in i(\mathfrak{b}_{0})\subset\mathfrak{b}\subset\mathfrak{b}\otimes A_{\mathfrak{g}}(O)$，其中 $\mathfrak{b}_{0}:=\operatorname{Lie}B_{0}$，而 $B_{0}\subset SL_{2}$ 是上三角矩阵群。这意味着，我们把 $\operatorname{Aut}_{2}^{0}O$-等变丛 $\mathfrak{F}^{0}_{B}$ 与 $\operatorname{Spec}\mathbb{C}$ 上某个 $\operatorname{Aut}_{2}^{0}O$-等变 $B$-丛的拉回等同起来，而后者来自某个态射 $\operatorname{Aut}_{2}^{0}O\to B_{0}\to B$（参看 3.1.10 的注 (iii)）。

\bdp{3.5.7} 设 $A$ 是带 $\operatorname{Aut}O$ 作用的代数。则 $\operatorname{Der}O$ 作用于 $A$，$L_{0}$ 在 $A$ 上的作用可对角化，且 $L_{0}:A\to A$ 的特征值为整数。假设 $L_{0}:A/\mathbb{C}\to A/\mathbb{C}$ 的特征值为正。则 $A=\mathbb{C}\oplus A_{+}$，其中 $A_{+}$ 是 $A$ 中对应于正特征值的所有特征子空间之和。$A_{+}$ 是 $A$ 的唯一 $L_{0}$-不变极大理想。$\operatorname{Spec}A$ 中对应的点记作 $0$。由于 $[L_{0},L_{n}]=-nL_{n}$，有 $L_{-1}A_{+}\subset A_{+}$。假设
\begin{equation}\tag{71}
L_{1}A_{+}\subset A_{+}.
\end{equation}
特别地，若 $L_{0}$ 在 $A/\mathbb{C}$ 上的特征值大于 $1$，则 (71) 成立；例如 $A=A_{\mathfrak{g}}(O)$ 时如此（见 (66) 与 (65)）。

假设 $G$ 是 $\mathfrak{g}$ 对应的伴随群。设 $\mathcal{E}$ 是 $\operatorname{Spec}A$ 上的 $\operatorname{Aut}O$-等变 $G$-丛。代数 $\mathbb{C}L_{-1}+\mathbb{C}L_{0}+\mathbb{C}L_{1}\simeq\mathfrak{sl}_{2}$ 稳定 $0\in\operatorname{Spec}A$，故它作用在 $\mathcal{E}$ 于 $0$ 处的纤维上。于是得到态射 $\sigma:\mathfrak{sl}_{2}\to\mathfrak{g}$，它在共轭意义下定义。

\SourcePage{85}
\emph{例.} 点 $0\in\operatorname{Spec}A_{\mathfrak{g}}(O))=\operatorname{Op}_{\mathfrak{g}}(O)$ 是 $\mathfrak{sl}_{2}$-oper 沿主嵌入 $\mathfrak{sl}_{2}\to\mathfrak{g}$ 的前推，该 $\mathfrak{sl}_{2}$-oper 对应于 Sturm--Liouville 算子 $(d/dt)^{2}$。若 $A=A_{\mathfrak{g}}(O)$ 且 $\mathcal{E}=\mathfrak{F}^{0}_{G}$，则 $\sigma$ 是主嵌入。
% [Tutor Note: 原书 p.85 此处“$\operatorname{Spec} A_{\mathfrak{g}}(O))$”多一个右括号，底本如此，保留原样。]

\bdp{3.5.8} \emph{命题.}
\begin{enumerate}[label=\arabic*),leftmargin=2.2em]
\item 以下条件等价：
\begin{enumerate}[label=\alph*),leftmargin=2.2em]
\item $\operatorname{Aut}O$-等变 $G$-丛 $\mathcal{E}$ 同构于 $\varphi^{*}\mathfrak{F}^{0}_{G}$，其中 $\varphi:\operatorname{Spec}A\to\operatorname{Op}_{\mathfrak{g}}(O)$ 是某个 $\operatorname{Aut}O$-等变态射；
\item $\mathcal{E}$ 上存在 $\operatorname{Aut}^{0}O$-不变的 $B$-结构，使得相应的 $\operatorname{Aut}^{0}O$-等变 $H$-丛同构于 3.5.5\footnote{由于 $G$ 是伴随群，$\operatorname{Aut}_{2}^{0}O$ 在 $F_{H}$ 上的作用经 $\operatorname{Aut}^{0}O$ 分解。} 所定义的、$\operatorname{Spec}\mathbb{C}$ 上 $\operatorname{Aut}^{0}O$-等变 $H$-丛 $F_{H}$ 的拉回；
\item $\sigma:\mathfrak{sl}_{2}\to\mathfrak{g}$ 是主嵌入。
\end{enumerate}
\item a) 中提到的态射 $\varphi$ 与同构 $\mathcal{E}\xrightarrow{\sim}\varphi^{*}\mathfrak{F}^{0}_{G}$ 都唯一。
\item b) 中提到的 $B$-结构唯一。
\end{enumerate}

\emph{证明.} 由 3.5.5，a) 蕴含 b)。为从 b) 推出 c)，只需看 $0\in\operatorname{Spec}A$ 上的情形。下面从 b) 推出 a)，并说明 2) 由 3) 推出。为此只需证明：若固定具有 b) 所述性质的 $B$-结构 $\mathcal{E}_{B}\subset\mathcal{E}$，则恰有一种方式构造 $\operatorname{Aut}O$-等变态射 $\varphi:\operatorname{Spec}A\to\operatorname{Op}_{\mathfrak{g}}(O)$ 与同构 $f:\mathcal{E}\xrightarrow{\sim}\varphi^{*}\mathfrak{F}^{0}_{G}$，使得 $f(\mathcal{E}_{B})=\varphi^{*}\mathfrak{F}^{0}_{B}$。由 3.5.3，$\mathcal{E}$ 与 $\mathcal{E}_{B}$ 给出 $\operatorname{Spec}(A\widehat\otimes O)$ 上的 $G$-丛 $\widetilde{\mathcal{E}}_{G}$，带 $B$-结构 $\widetilde{\mathcal{E}}_{B}\subset\widetilde{\mathcal{E}}_{G}$、沿 $\operatorname{Spec}O$ 的联络 $\nabla$，以及 $\operatorname{Aut}O$ 在 $(\widetilde{\mathcal{E}}_{G},\widetilde{\mathcal{E}}_{B},\nabla)$ 上的作用。现在 $\varphi$ 与 $f$ 的唯一性清楚；为证其存在，须证 $(\widetilde{\mathcal{E}}_{G},\widetilde{\mathcal{E}}_{B},\nabla)$ 是 oper 族，即须证 3.1.2 定义的 $c(\nabla)$ 满足定义 3.1.3 的条件 1 与 2。在我们情形，$c(\nabla)$ 是 $(\mathfrak{g}/\mathfrak{b})_{\widetilde{\mathcal{E}}}\otimes_{O}\omega_{O}$ 的 $\operatorname{Aut}O$-不变截面；只需对其到 $\operatorname{Spec}A\subset\operatorname{Spec}(A\widehat\otimes O)$ 的限制 $c_{0}(\nabla)$ 验证条件 1 与 2。

\SourcePage{86}
$c_{0}(\nabla)$ 是 $H^{0}(\operatorname{Spec}A,(\mathfrak{g}/\mathfrak{b})_{\mathcal{E}_{B}})\otimes\omega_{O}/t\omega_{O}$ 的 $\operatorname{Aut}^{0}O$-不变元。令 $\operatorname{gr}^{k}\mathfrak{g}$ 与 3.1.1 中含义相同。既然知道 $\mathcal{E}_{B}$ 对应的 $H$-丛，便看出存在 $\operatorname{Aut}^{0}O$-等变同构
\begin{equation}\tag{72}
H^{0}(\operatorname{Spec}A,\operatorname{gr}^{k}\mathfrak{g}_{\mathcal{E}_{B}})\otimes\omega_{O}/t\omega_{O}\xrightarrow{\sim}A\otimes(\omega_{O}/t\omega_{O})^{\otimes(k+1)}\otimes\operatorname{gr}^{k}\mathfrak{g}.
\end{equation}
由于 $L_{0}$ 在 $(\omega_{O}/t\omega_{O})^{\otimes(k+1)}$ 上按乘以 $-k-1$ 的方式作用，$A\otimes(\omega_{O}/t\omega_{O})^{\otimes(k+1)}$ 的 $\operatorname{Aut}^{0}O$-不变部分在 $k<-1$ 时为 $0$，在 $k=-1$ 时为 $\mathbb{C}$。因此
\[
c_{0}(\nabla)\in\operatorname{gr}^{-1}\mathfrak{g}\subset A\otimes\operatorname{gr}^{-1}\mathfrak{g}=H^{0}(\operatorname{Spec}A,\operatorname{gr}^{-1}\mathfrak{g}_{\mathcal{E}_{B}})\otimes\omega_{O}/t\omega_{O}.
\]
这样我们验证了 3.1.3 的条件 1；还需在 $\operatorname{Spec}A$ 的某一点上验证条件 2，例如在 $0\in\operatorname{Spec}A$ 上。以 $(\widetilde{\mathcal{E}}^{0}_{G},\widetilde{\mathcal{E}}^{0}_{B},\nabla)$ 记 $(\widetilde{\mathcal{E}}_{G},\widetilde{\mathcal{E}}_{B},\nabla)$ 到 $\{0\}\times\operatorname{Spec}O\subset\operatorname{Spec}(A\widehat\otimes O)$ 的限制。则 $\widetilde{\mathcal{E}}^{0}_{G}$ 是平凡 $G$-丛，$\nabla$ 是平凡联络，$\mathfrak{sl}_{2}$ 经 3.5.7 提到的态射 $\sigma:\mathfrak{sl}_{2}\to\mathfrak{g}$ 与嵌入 $\mathfrak{sl}_{2}=\mathbb{C}L_{-1}+\mathbb{C}L_{0}+\mathbb{C}L_{1}\hookrightarrow\operatorname{Der}O$ 作用在 $(\widetilde{\mathcal{E}}^{0}_{G},\nabla)$ 上，且 $\widetilde{\mathcal{E}}^{0}_{B}$ 关于 $\mathfrak{sl}_{2}$ 不变。由于 $\sigma$ 是主嵌入，$(\widetilde{\mathcal{E}}^{0}_{G},\widetilde{\mathcal{E}}^{0}_{B},\nabla)$ 是 $0\in\operatorname{Op}_{\mathfrak{g}}(O)$ 对应的 oper。

下面证 3)。置 $\mathfrak{a}=H^{0}(\operatorname{Spec}A,\mathfrak{g}_{\mathcal{E}})$，$\mathfrak{a}_{k}:=\{a\in\mathfrak{a}\mid L_{0}a=ka\}$。若固定 $\mathcal{E}$ 上的 $B$-结构，则 3.1.1 的滤链 $\mathfrak{g}^{k}$ 在 $\mathfrak{a}$ 上诱导滤链 $\mathfrak{a}^{k}$。若该 $B$-结构具有 b) 所述性质，则 $\mathfrak{a}^{k}$ 是 $\operatorname{Aut}^{0}O$-不变的，且 $\mathfrak{a}^{k}/\mathfrak{a}^{k+1}$ 作为 $\operatorname{Aut}^{0}O$-模同构于 $A\otimes(\omega_{O}/t\omega_{O})^{\otimes k}\otimes\operatorname{gr}^{k}\mathfrak{g}$（见 (72)）。因此 $L_{0}$ 在 $\mathfrak{a}^{k}/\mathfrak{a}^{k+1}$ 上的特征值 $\ge-k$，且 $A$-模 $\mathfrak{a}^{k}/\mathfrak{a}^{k+1}$ 由其对应于特征值 $-k$ 的 $L_{0}$-特征向量生成。故
\begin{equation}\tag{73}
\mathfrak{a}^{k}=\sum_{i\le-k}A\mathfrak{a}_{i}.
\end{equation}
$\mathcal{E}$ 上的 $B$-结构可由 Borel 子代数 $\mathfrak{a}^{0}\subset\mathfrak{a}$ 重构。

剩下从 c) 推出 b)。按 (73) 定义 $\mathfrak{a}^{k}$。由于 $\mathfrak{a}$ 是有限型的自由 $L_{0}$-分次 $A$-模，$\mathfrak{a}^{k}/\mathfrak{a}^{k+1}$ 也是。因此 $\mathfrak{a}^{k}$ 定义 $\mathfrak{g}_{\mathcal{E}}$ 的一个向量\SourcePage{87}%
子丛。若 $k=0$，该子丛是李子代数，故它定义截面 $s:\operatorname{Spec}A\to S_{\mathcal{E}}$，其中 $S$ 是 $\mathfrak{g}$ 的子代数概形。一个无穷小计算表明，态射 $G/B\to S$，$g\mapsto g\mathfrak{b}g^{-1}$，是开嵌入；又因 $G/B$ 射影，它也是闭嵌入。由 c)，$s(0)\in(G/B)_{\mathcal{E}}\subset S_{\mathcal{E}}$，故 $s(\operatorname{Spec}A)\subset(G/B)_{\mathcal{E}}$，于是 $s$ 定义 $\mathcal{E}$ 上的 $B$-结构。显然它是 $\operatorname{Aut}^{0}O$-不变的。$\operatorname{Spec}A$ 上相应的 $\operatorname{Aut}^{0}O$-等变 $H$-丛是 $\operatorname{Spec}\mathbb{C}$ 上某个 $\operatorname{Aut}^{0}O$-等变 $H$-丛 $F$ 的拉回（对 $\operatorname{Spec}A$ 上任意 $\operatorname{Aut}^{0}O$-等变 $H$-丛与任意环面 $H$ 都成立；事实上可设 $H=\mathbb{G}_{m}$，把 $\mathbb{G}_{m}$-丛解释为线丛，并利用有限生成分次投射 $A$-模自由这一事实）。为求 $F$，看 $0\in\operatorname{Spec}A$ 上的情形。

\emph{注.} 命题 3.5.8 的证明表明：若 c) 成立，则 $\mathcal{E}$ 上存在唯一的 $\operatorname{Aut}^{0}O$-不变 $B$-结构。

\bdp{3.5.9} \emph{推论.} 若 $G$ 是伴随群，则偶 $(\operatorname{Op}_{\mathfrak{g}}(O),\mathfrak{F}^{0}_{G})$ 没有非平凡的 $\operatorname{Aut}O$-等变自同构。
这是命题 3.5.8 的陈述 2) 取 $A=A_{\mathfrak{g}}(O)$ 的情形。

\bdp{3.5.10} 回忆：交换 $\mathbb{C}$-代数 $R$ 上的李代数胚，是李 $\mathbb{C}$-代数 $\mathfrak{a}$ 配以 $R$-模结构及映射 $\varphi:\mathfrak{a}\to\operatorname{Der}R$，使得 1) $\varphi$ 是李代数态射与 $R$-模态射；2) 对 $a_{1},a_{2}\in\mathfrak{a}$ 与 $f\in R$，有 $[a_{1},fa_{2}]=f[a_{1},a_{2}]+v(f)a_{2}$，其中 $v:=\varphi(a_{1})$。

\emph{注.}
(i) [Ma87] 与 [Ma96] 是李代数胚与李群胚的标准文献。另见 [We] 与 [BB93]。本文只需李代数胚的定义。
(ii) 李代数胚也称为 $(\mathbb{C},R)$-李代数（见 [R]），并有各种其他名称（见 [Ma96]）。

\bdp{3.5.11} 以 $\mathfrak{a}_{\mathfrak{g}}$ 记 $\mathfrak{F}^{0}_{G}$ 的（整体）无穷小对称空间。$\mathfrak{a}_{\mathfrak{g}}$ 的元素是偶，由 $\operatorname{Op}_{\mathfrak{g}}(O)=\operatorname{Spec}A_{\mathfrak{g}}(O)$ 上的向量场\SourcePage{88}%
（即 $A_{\mathfrak{g}}(O)$ 的导子）及其到主 $G$-丛 $\mathfrak{F}^{0}_{G}$ 上 $G$-不变向量场的提升组成。$\mathfrak{a}_{\mathfrak{g}}$ 是 $A_{\mathfrak{g}}(O)$ 上的李代数胚。我们有典范正合列
\[
0\to\mathfrak{g}_{\mathrm{univ}}\to\mathfrak{a}_{\mathfrak{g}}\to\operatorname{Der}A_{\mathfrak{g}}(O)\to0,
\]
其中 $\mathfrak{g}_{\mathrm{univ}}$ 是 $\mathfrak{g}$ 的 $\mathfrak{F}^{0}_{G}$-扭的整体截面空间。当然，把 $G$ 换成伴随群 $G_{ad}$ 时，$\mathfrak{a}_{\mathfrak{g}}$ 与 $\mathfrak{g}_{\mathrm{univ}}$ 不变。故 $\mathfrak{a}_{\mathfrak{g}}$ 与 $\mathfrak{g}_{\mathrm{univ}}$ 不依赖 $\omega_{O}^{1/2}$ 的选取。

$\operatorname{Der}O$ 在 $\mathfrak{F}^{0}_{G}$ 上的作用诱导李代数态射 $\operatorname{Der}O\to\mathfrak{a}_{\mathfrak{g}}$。特别地，$\operatorname{Der}O$ 作用于 $\mathfrak{a}_{\mathfrak{g}}$。

\bdp{3.5.12} \emph{引理.} $\mathfrak{a}_{\mathfrak{g}}$ 在 $\mathfrak{g}_{\mathrm{univ}}$ 上的伴随表示定义了 $\mathfrak{a}_{\mathfrak{g}}$ 与 $\mathfrak{g}_{\mathrm{univ}}$ 的无穷小对称代数胚之间的同构。

\bdp{3.5.13} \emph{命题.} 偶 $(A_{\mathfrak{g}}(O),\mathfrak{a}_{\mathfrak{g}})$ 的 $\operatorname{Der}O$-等变自同构群等于 $\operatorname{Aut}\Gamma$，其中 $\Gamma$ 是 $\mathfrak{g}$ 的 Dynkin 图。

\emph{证明.} 设 $G$ 是 $\mathfrak{g}$ 对应的伴随群。以 $L$ 记 $(A_{\mathfrak{g}}(O),\mathfrak{g}_{\mathrm{univ}})$ 的 $\operatorname{Der}O$-等变自同构群。由 3.5.12，须证 $L=\operatorname{Aut}\Gamma$。有显然态射 $i:\operatorname{Aut}\Gamma=\operatorname{Aut}(G,B)/B\to L$ 与 $\pi:L\to\operatorname{Aut}\Gamma$，使得 $\pi i=\mathrm{id}$。$\operatorname{Ker}\pi$ 是 $(\operatorname{Op}_{\mathfrak{g}}(O),\mathfrak{F}^{0}_{G})$ 的 $\operatorname{Der}O$-等变自同构群；故由 3.5.9，$\operatorname{Ker}\pi$ 平凡。\hfill$\Box$

\bdp{3.5.14} \emph{命题.} 偶 $(A_{\mathfrak{g}}(O),\mathfrak{a}_{\mathfrak{g}})$ 没有能诱导 $\operatorname{gr}A_{\mathfrak{g}}(O)$ 平凡自同构的非平凡 $\operatorname{Der}O$-等变自同构（$\operatorname{gr}$ 对应 3.2.1 的滤链）。

\emph{证明.} 设 $\Gamma$ 是 $\mathfrak{g}$ 的 Dynkin 图。由 3.5.13 与 (48)，须证 $\operatorname{Aut}\Gamma$ 在 2.7.1 的代数 $\mathfrak{z}^{cl}_{{}^{L}\mathfrak{g}}(O)$ 上的作用忠实。故只需证 $\operatorname{Aut}\Gamma$ 在 $W\backslash\mathfrak{h}$ 上的作用忠实（$W$ 表示 Weyl 群）。设 $C\subset\operatorname{Aut}\mathfrak{h}$ 是\SourcePage{89}%
根系的自同构群。存在 $a\in\mathfrak{h}$，其在 $C$ 中的稳定化子平凡。于是 $\operatorname{Aut}\Gamma=C/W$ 在 $W\backslash\mathfrak{h}$ 上的作用忠实。\hfill$\Box$

\bdp{3.5.15} 给 $\mathfrak{a}_{\mathfrak{g}}$ 配备最弱的平移不变拓扑，使得 $\mathfrak{F}^{0}_{G}$ 全空间上任一正则函数的稳定化子为开集（回忆 $\mathfrak{a}_{\mathfrak{g}}$ 作用于 $\mathfrak{F}^{0}_{G}$）。这是使 $\mathfrak{g}_{\mathrm{univ}}$ 每个元素的 $\mathfrak{a}_{\mathfrak{g}}$-中心化子为开集的最弱平移不变拓扑。于是该拓扑可由 $\mathfrak{a}_{\mathfrak{g}}$ 上的李代数胚结构重构。显然典范态射 $\operatorname{Der}O\to\mathfrak{a}_{\mathfrak{g}}$ 连续。

\bdp{3.5.16} 以 $\mathfrak{a}_{\mathfrak{b}}$ 记 $\mathfrak{F}^{0}_{B}$ 的（整体）无穷小对称李代数胚。以 $\mathfrak{b}_{\mathrm{univ}}$（相应地 $\mathfrak{n}_{\mathrm{univ}}$）记 $\mathfrak{b}$（相应地 $\mathfrak{n}$）的 $\mathfrak{F}^{0}_{B}$-扭的整体截面空间。有典范正合列
\[
0\to\mathfrak{b}_{\mathrm{univ}}\to\mathfrak{a}_{\mathfrak{b}}\to\operatorname{Der}A_{\mathfrak{g}}(O)\to0.
\]
$\mathfrak{a}_{\mathfrak{b}}$ 是 $\mathfrak{a}_{\mathfrak{g}}$ 的子代数胚；事实上 $\mathfrak{a}_{\mathfrak{b}}$ 是 $\mathfrak{b}_{\mathrm{univ}}\subset\mathfrak{a}_{\mathfrak{g}}$ 的正规化子。$\operatorname{Der}^{0}O$ 在 $\mathfrak{a}_{\mathfrak{g}}$ 中的像含于 $\mathfrak{a}_{\mathfrak{b}}$。

$\mathfrak{n}_{\mathrm{univ}}$ 是 $\mathfrak{a}_{\mathfrak{b}}$ 的理想，$\mathfrak{a}_{\mathfrak{b}}/\mathfrak{n}_{\mathrm{univ}}$ 是 $\mathfrak{F}^{0}_{H}$ 的（整体）无穷小对称代数胚。由于 $\mathfrak{F}^{0}_{H}$ 平凡且其平凡化「几乎」唯一（见 3.5.5），$\mathfrak{a}_{\mathfrak{b}}/\mathfrak{n}_{\mathrm{univ}}$ 典范同构于 $\operatorname{Der}A_{\mathfrak{g}}(O)$ 与 $A_{\mathfrak{g}}(O)\otimes\mathfrak{h}$ 的半直和。以 $\mathfrak{a}_{\mathfrak{n}}$ 记 $\operatorname{Der}A_{\mathfrak{g}}(O)\subset\mathfrak{a}_{\mathfrak{b}}/\mathfrak{n}_{\mathrm{univ}}$ 在 $\mathfrak{a}_{\mathfrak{b}}$ 中的原像。

\emph{注.} 由 3.5.5，复合 $\operatorname{Der}^{0}O\to\mathfrak{a}_{\mathfrak{b}}/\mathfrak{n}_{\mathrm{univ}}=\operatorname{Der}A_{\mathfrak{g}}(O)\oplus(A_{\mathfrak{g}}(O)\otimes\mathfrak{h})$ 的像含于 $\operatorname{Der}A_{\mathfrak{g}}(O)\oplus\mathfrak{h}$；它等于自然态射 $\operatorname{Der}^{0}O\to\operatorname{Der}A_{\mathfrak{g}}(O)$ 与态射 $\operatorname{Der}^{0}O\to\mathfrak{h}$（$L_{0}\mapsto-\check\rho$，$n>0$ 时 $L_{n}\mapsto0$）之和。

\bdp{3.5.17} 下面用 $L_{0}$ 在 $\mathfrak{a}_{\mathfrak{g}}$ 上的作用描述 $\mathfrak{a}_{\mathfrak{b}}$、$\mathfrak{b}_{\mathrm{univ}}$ 等。将使用如下记号。若 $\operatorname{Der}O$ 作用在拓扑向量空间 $V$ 上，使 $L_{0}:V\to V$ 的特征值为整数，则以 $V^{\le k}$ 记 $V$ 中含所有满足 $L_{0}v=nv$、$n\le k$ 的 $v$ 的最小闭子空间。置 $V^{<k}:=V^{\le k-1}$。若 $V$ 是某代数 $A$ 上的拓扑模，\SourcePage{90}%
$W$ 是 $V$ 的子空间，则以 $A\cdot W$ 记 $V$ 中含所有 $aw$（$a\in A$，$w\in W$）的最小闭子空间。

\bdp{3.5.18} \emph{命题.} i) 下列等式成立：
\begin{equation}\tag{74}
\mathfrak{b}_{\mathrm{univ}}=A_{\mathfrak{g}}(O)\cdot(\mathfrak{g}_{\mathrm{univ}})^{\le0}
\end{equation}
\begin{equation}\tag{75}
\mathfrak{n}_{\mathrm{univ}}=A_{\mathfrak{g}}(O)\cdot\mathfrak{g}_{\mathrm{univ}}^{<0}
\end{equation}
\begin{equation}\tag{76}
\mathfrak{a}_{\mathfrak{b}}=A_{\mathfrak{g}}(O)\cdot(\mathfrak{a}_{\mathfrak{g}})^{\le0}
\end{equation}
\begin{equation}\tag{77}
\mathfrak{a}_{\mathfrak{n}}=A_{\mathfrak{g}}(O)\cdot\mathfrak{a}_{\mathfrak{g}}^{<0}
\end{equation}

ii) 态射
\[
\begin{aligned}
(\mathfrak{a}_{\mathfrak{g}})^{\le0}&\longrightarrow
A_{\mathfrak{g}}(O)(\mathfrak{a}_{\mathfrak{g}})^{\le0}/A_{\mathfrak{g}}(O)\mathfrak{a}_{\mathfrak{g}}^{<0}\\
&=\mathfrak{a}_{\mathfrak{b}}/\mathfrak{a}_{\mathfrak{n}}=A_{\mathfrak{g}}(O)\otimes\mathfrak{h}
\end{aligned}
\]
的像等于 $\mathfrak{h}$，故有典范同构
\begin{equation}\tag{78}
(\mathfrak{a}_{\mathfrak{g}})^{\le0}\big/\bigl(A_{\mathfrak{g}}(O)\cdot\mathfrak{a}_{\mathfrak{g}}^{<0}\cap(\mathfrak{a}_{\mathfrak{g}})^{\le0}\bigr)\xrightarrow{\sim}\mathfrak{h}.
\end{equation}

\emph{证明.} i) (74)–(77) 由 (69) 推出。或者注意到 (74) 与 (75) 是 (73) 的特例，并如下证明 (76)。由 (74)，$A_{\mathfrak{g}}(O)\cdot(\mathfrak{a}_{\mathfrak{g}})^{\le0}\supset\mathfrak{b}_{\mathrm{univ}}$，且 $A_{\mathfrak{g}}(O)\cdot(\operatorname{Der}A_{\mathfrak{g}}(O))^{\le0}=\operatorname{Der}A_{\mathfrak{g}}(O)$，故 $A_{\mathfrak{g}}(O)\cdot(\mathfrak{a}_{\mathfrak{g}})^{\le0}\supset\mathfrak{a}_{\mathfrak{b}}$。又由 (74)，$(\mathfrak{a}_{\mathfrak{g}}/\mathfrak{a}_{\mathfrak{b}})^{\le0}=(\mathfrak{g}_{\mathrm{univ}}/\mathfrak{b}_{\mathrm{univ}})^{\le0}=(\mathfrak{g}_{\mathrm{univ}})^{\le0}/(\mathfrak{b}_{\mathrm{univ}})^{\le0}=0$，故 $A_{\mathfrak{g}}(O)\cdot(\mathfrak{a}_{\mathfrak{g}})^{\le0}\subset\mathfrak{a}_{\mathfrak{b}}$。
ii) $(\mathfrak{a}_{\mathfrak{g}})^{\le0}$ 在 $A_{\mathfrak{g}}(O)\otimes\mathfrak{h}$ 中的像等于 $(A_{\mathfrak{g}}(O)\otimes\mathfrak{h})^{\le0}=\mathfrak{h}$。\hfill$\Box$

\bdsec{3.6. Feigin--Frenkel 同构 II}
\bdp{3.6.1} 设 $A$ 是 $\mathbb{C}[h]$ 上的结合代数，且作为 $\mathbb{C}[h]$-模平坦。置 $A_{0}:=A/hA$。以 $\mathfrak{z}$ 记 $A_{0}$ 的中心。若 $\mathfrak{z}=A_{0}$，即 $A_{0}$ 交换，则 $\mathfrak{z}$ 带有标准泊松括号
\begin{equation}\tag{79}
\{z_{1},z_{2}\}:=[\tilde z_{1},\tilde z_{2}]/h\bmod h.
\end{equation}
